\subsection{Coarse-to-Fine Cascade}\label{sec:results-cascade}

This section evaluates the Cascade configuration (\autoref{tab:configs}): how
accuracy develops across levels, how well the coarse level localizes the
target, how the cascade compares with Medverse in accuracy and latency, and
which prior works better during training.

\paragraph{Accuracy across levels.}
\autoref{tab:cascade-levels} reports the Dice of each level, with each crop
placed from the previous level's prediction. Dice increases from level to
level on every dataset, e.g.\ from 0.603 to 0.751 on FLARE22 and from 0.065
to 0.152 on HU\_LWK1. The stitched mask keeps the coarser predictions outside
the finest crop. This recovers target parts on TotalSeg MRI (0.495 to 0.503),
adds false positives on MSD Prostate (0.381 to 0.375), and changes nothing
where the coarser levels predict no foreground outside the finest crop.

\begin{table}[t]
  \centering
  \small
  \caption{Dice of the Cascade model at each level, with crops placed using the
  previous level's prediction, and Dice of the final stitched prediction.}
  \label{tab:cascade-levels}
  \begin{tabular}{@{}lccc@{}}
    \toprule
    Dataset & Spacings (mm) & Dice per level & Stitched \\
    \midrule
    % TotalSeg CT  & $[6, 3, 1.5]$    & $0.394 \to 0.522 \to 0.576$ & 0.585 \\ % removed for final paper (see TODO.md)
    FLARE22      & $[6, 3, 2.5]$    & $0.603 \to 0.718 \to 0.751$ & 0.751 \\
    MSD Prostate & $[3, 1.5, 0.75]$ & $0.276 \to 0.361 \to 0.381$ & 0.375 \\
    HU\_LWK1     & $[6, 3, 1]$      & $0.065 \to 0.109 \to 0.152$ & 0.152 \\
    TotalSeg MRI & $[6, 3, 1.5]$    & $0.365 \to 0.479 \to 0.495$ & 0.503 \\
    NasalSeg     & $[1.5, 0.8]$     & $0.484 \to 0.547$           & 0.547 \\
    \bottomrule
  \end{tabular}
\end{table}

\autoref{fig:c2f-hu-lwk1} and \autoref{fig:c2f-flare22} show the crop
placement for Cascade and for Medverse run through the same cascade.

\begin{figure}[t]
  \centering
  \includegraphics[width=\linewidth]{c2f_hu_lwk1.png}
  \caption{HU\_LWK1: Cascade's box tracks the vertebra center across all
  three levels; Medverse's drifts off by level 1, so its finest crop no
  longer contains the target.}
  \label{fig:c2f-hu-lwk1}
\end{figure}

\begin{figure}[t]
  \centering
  \includegraphics[width=\linewidth]{c2f_flare22.png}
  \caption{FLARE22: both models' boxes converge on the spleen by level 1.}
  \label{fig:c2f-flare22}
\end{figure}


\paragraph{Localization.}
On TotalSeg CT, a middle-level crop centered at the centroid of the coarse
prediction contains on average 92.2\% of the target's voxels, compared with
97.0\% when centered at the ground-truth centroid. Localization thus loses
about 5\% of the target at this transition, and most of the remaining error
arises within the fine crops. The middle-to-finest transition is not measured.

\paragraph{Comparison with Medverse.}
\autoref{tab:cascade-medverse} compares Cascade with Medverse's released
weights, run with its own autoregressive inference and with its number of
steps matched to our number of levels per dataset; its native depth is
compared below. MSD Hippocampus runs at a single level, so its row compares
the two models without a cascade. Medverse is not evaluated on TotalSeg CT.

\begin{table}[t]
  \centering
  \small
  \caption{Cascade vs.\ Medverse: macro Dice and mean latency per sample.
  Medverse is evaluated with its number of autoregressive steps matched to our
  number of levels per dataset.}
  \label{tab:cascade-medverse}
  \begin{tabular}{@{}lcccccc@{}}
    \toprule
    & & \multicolumn{2}{c}{Dice} & & \multicolumn{2}{c}{Latency (ms)} \\
    \cmidrule(lr){3-4} \cmidrule(lr){6-7}
    Dataset & Spacings (mm) & Cascade & Medverse & & Cascade & Medverse \\
    \midrule
    FLARE22         & $[6, 3, 2.5]$    & \textbf{0.751} & 0.359 & & 311 & 4894 \\
    MSD Prostate    & $[3, 1.5, 0.75]$ & \textbf{0.375} & 0.357 & & 279 & 5006 \\
    MSD Hippocampus & $[0.5]$          & 0.303 & \textbf{0.686} & & 99  & 50   \\
    HU\_LWK1        & $[6, 3, 1]$      & \textbf{0.152} & 0.020 & & 439 & 5336 \\
    TotalSeg MRI    & $[6, 3, 1.5]$    & \textbf{0.503} & 0.055 & & 194 & 5188 \\
    NasalSeg        & $[1.5, 0.8]$     & 0.547 & \textbf{0.737} & & 161 & 1886 \\
    \midrule
    Mean (sample-weighted)$^{\dagger}$ & & \textbf{0.526} & 0.393 & & \textbf{206} & 3491 \\
    \midrule
    ISLES22         & $[3, 1.5]$       & \textbf{0.069} & 0.059 & & 155 & 1894 \\
    Shifts-MS       & $[3, 1.5]$       & 0.031 & \textbf{0.098} & & 289 & 1911 \\
    ATLAS v2.0      & $[3.8, 1.9]$     & \textbf{0.026} & 0.024 & & 149 & 1862 \\
    GNC\_705        & $[6, 3, 1.2]$    & \textbf{0.024} & 0.007 & & 288 & 4874 \\
    \bottomrule
  \end{tabular}
  \par\vspace{0.5em}
  \raggedright\footnotesize
  $^{\dagger}$Over the six datasets above the second midrule, each weighted by
  its number of evaluated tasks (36 for HU\_LWK1 to 650 for FLARE22).
  Unweighted: Dice 0.439 vs.\ 0.369, latency 247 vs.\ 3727\,ms. The four
  lesion datasets are excluded, since both models score below 0.1 on all of
  them.
\end{table}

On the six non-lesion datasets, Cascade is more accurate on four (FLARE22,
MSD Prostate, HU\_LWK1, TotalSeg MRI) and Medverse on two (MSD Hippocampus,
NasalSeg); the sample-weighted means are 0.526 and 0.393. The largest gaps in
our favor are on TotalSeg MRI, which is part of our training data, and on
FLARE22, whose classes are mostly in our training vocabulary. Medverse's two
wins, in contrast, are on datasets that neither model saw during training. On
the four lesion datasets, both models stay below 0.1 Dice.

\autoref{fig:cascade-medverse-qual} and \autoref{fig:cascade-medverse-qual-fail}
show one case from each side, with Medverse at its native depth. Ground truth
is shown in green, Cascade in blue, and Medverse in orange. On a FLARE22
gallbladder, Cascade recovers the organ (Dice 0.79), while Medverse predicts a
small, mislocated fragment (0.00). On MSD Hippocampus, Cascade's prediction
breaks into disconnected blobs (0.34), while Medverse's is a single region
close to the ground truth (0.70).

\begin{figure}[t]
  \centering
  \includegraphics[width=0.85\linewidth]{flare22_gallbladder_win.png}
  \caption{FLARE22 gallbladder, subject FLARE22\_Tr\_0001: Cascade recovers
  the organ, while Medverse (native depth) misses it almost entirely.}
  \label{fig:cascade-medverse-qual}
\end{figure}

\begin{figure}[t]
  \centering
  \includegraphics[width=0.85\linewidth]{msd_hippocampus_fail.png}
  \caption{MSD Hippocampus, posterior class, subject hippocampus\_001:
  Cascade's prediction breaks into disconnected blobs, while Medverse (native
  depth) predicts a single contiguous region.}
  \label{fig:cascade-medverse-qual-fail}
\end{figure}
% Provenance: re-rendered from run_cascade / model.predict() on these
% (target, context) pairs. FLARE22: 2026-09-26_147_cascade_randomfg_predprior_regon_flare22_ood
% vs. 2026-09-15_medverse_ar_flare22. Hippocampus:
% 2026-09-24_146_cascade_randomfg_predprior_regoff_msd_hippocampus_ood vs.
% 2026-09-14_medverse_ar_msd_hippocampus. First sample the eval loader hit
% per class, not a best/worst search.

Cascade is about $17\times$ faster than step-matched Medverse (206 vs.\
3491\,ms per sample, sample-weighted; $15\times$ unweighted, 247 vs.\
3727\,ms). At a single level (MSD Hippocampus), Medverse is faster (50 vs.\
99\,ms). The speedup therefore comes from processing one crop per level
instead of sliding a window over the whole volume at each step, not from a
cheaper forward pass.

\paragraph{Medverse at native depth.}
By default, Medverse runs more autoregressive steps than our cascade has
levels. \autoref{tab:medverse-matched-native} compares its Dice at both
depths.

\begin{table}[t]
  \centering
  \small
  \caption{Medverse's macro Dice with its number of autoregressive steps
  matched to our number of levels (matched, as in
  \autoref{tab:cascade-medverse}) and at its default number of steps
  (native).}
  \label{tab:medverse-matched-native}
  \begin{tabular}{@{}lccc@{}}
    \toprule
    Dataset & Spacings (mm) & Matched & Native \\
    \midrule
    FLARE22         & $[6, 3, 2.5]$    & 0.359 & \textbf{0.440} \\
    MSD Prostate    & $[3, 1.5, 0.75]$ & 0.357 & \textbf{0.423} \\
    MSD Hippocampus & $[0.5]$          & 0.686 & \textbf{0.691} \\
    HU\_LWK1        & $[6, 3, 1]$      & 0.020 & \textbf{0.071} \\
    TotalSeg MRI    & $[6, 3, 1.5]$    & 0.055 & \textbf{0.072} \\
    NasalSeg        & $[1.5, 0.8]$     & \textbf{0.737} & 0.735 \\
    \midrule
    ISLES22         & $[3, 1.5]$       & \textbf{0.059} & 0.053 \\
    Shifts-MS       & $[3, 1.5]$       & 0.098 & \textbf{0.152} \\
    ATLAS v2.0      & $[3.8, 1.9]$     & \textbf{0.024} & 0.023 \\
    GNC\_705        & $[6, 3, 1.2]$    & 0.007 & \textbf{0.011} \\
    \bottomrule
  \end{tabular}
\end{table}

Matching lowers Medverse's Dice on seven of the ten datasets, most on FLARE22
($-0.081$) and MSD Prostate ($-0.066$), and raises it on the other three by at
most 0.006. At native depth, Medverse overtakes Cascade on MSD Prostate (0.423
vs.\ 0.375), so the two models split the non-lesion datasets three to three,
and it exceeds 0.1 on one lesion dataset, Shifts-MS (0.152). The accuracy
comparison therefore depends on how Medverse is run. The latency advantage
does not, since native depth only adds steps.

\paragraph{Medverse with our cascade.}
We also ran Medverse's weights inside our cascade inference
(\autoref{sec:methodology-cascade}) on seven held-out datasets, with and
without its coarse prediction passed on as a prior
(\autoref{tab:medverse-prior-appendix}). Here Medverse runs a multi-level
cascade on every dataset, including MSD Hippocampus ($[6, 3, 0.5]$\,mm), so
these numbers are not directly comparable with \autoref{tab:cascade-medverse}.
The prior lowers Dice on six of the seven datasets, most on MSD Hippocampus
(0.691 to 0.050) and ATLAS v2.0 (0.126 to 0.020); on HU\_LWK1, Dice is 0 in
both cases (\autoref{fig:c2f-hu-lwk1}). Medverse is trained with a downsampled
ground-truth mask as prior, so its own coarse errors are out of distribution
for this input.

\paragraph{Source of the label prior.}
\autoref{tab:cascade-ablation} compares two models that differ only in the
prior used during training: the model's own prediction from the previous
level, or a perturbed ground-truth mask. Both use their own prediction at
inference. The predicted prior is better on the three datasets from the
training vocabulary; on TotalSeg CT, it is better on 99 of 117 classes
(median $+0.043$) and raises NSD from 0.576 to 0.631. On the seven datasets
outside the training vocabulary, it is better on MSD Prostate, HU\_LWK1, and
NasalSeg, by 0.053 to 0.148. On the four lesion datasets, both priors stay
below 0.07 Dice and differ by at most 0.004.

\begin{table}[t]
  \centering
  \small
  \caption{Macro Dice of Cascade models trained with a perturbed ground-truth
  prior or with their own predicted prior. Both use their own prediction as
  prior at inference.}
  \label{tab:cascade-ablation}
  \begin{tabular}{@{}lcc@{}}
    \toprule
    Dataset & Ground-truth prior & Predicted prior \\
    \midrule
    TotalSeg CT     & 0.543 & \textbf{0.585} \\
    TotalSeg MRI    & 0.449 & \textbf{0.503} \\
    FLARE22         & 0.731 & \textbf{0.751} \\
    \midrule
    MSD Prostate    & 0.310 & \textbf{0.375} \\
    HU\_LWK1        & 0.099 & \textbf{0.152} \\
    NasalSeg        & 0.399 & \textbf{0.547} \\
    \midrule
    ISLES22         & 0.067 & \textbf{0.069} \\
    Shifts-MS       & \textbf{0.035} & 0.031 \\
    ATLAS v2.0      & \textbf{0.029} & 0.026 \\
    GNC\_705        & \textbf{0.027} & 0.024 \\
    \bottomrule
  \end{tabular}
\end{table}
